\section*{Appendix}

\subsection{Compliance and reproducibility}

All submitted panels are produced by fixed automated pipelines without manual pixel editing or per-story prompt rewrites after viewing outputs.
General rules---subject-type normalization for humans, animals, and robots---apply uniformly and are not story-specific hard coding.
The LLM is used only for text-side structured planning; image synthesis runs locally on open checkpoints (SDXL-Turbo, IP-Adapter, StoryDiffusion, PixArt-$\alpha$, CLIP).
The system is not an agent that inspects images and autonomously revises prompts; generation, optional scoring, routing, and logging are predetermined stages.

\subsection{Anchor bank selection}

\begin{figure}[!h]
  \centering
  \begin{subfigure}{0.22\columnwidth}
    \fbox{\parbox[c][0.7in][c]{0.95\linewidth}{\centering\tiny C1}}
    \caption{Cand.~1}
  \end{subfigure}
  \begin{subfigure}{0.22\columnwidth}
    \fbox{\parbox[c][0.7in][c]{0.95\linewidth}{\centering\tiny C2}}
    \caption{Cand.~2}
  \end{subfigure}
  \begin{subfigure}{0.22\columnwidth}
    \fbox{\parbox[c][0.7in][c]{0.95\linewidth}{\centering\tiny C3}}
    \caption{Cand.~3}
  \end{subfigure}
  \begin{subfigure}{0.22\columnwidth}
    \fbox{\parbox[c][0.7in][c]{0.95\linewidth}{\centering\tiny OK}}
    \caption{Canonical}
  \end{subfigure}
  \caption{Anchor bank candidates and selected canonical half-body anchor.}
  \label{fig:anchor}
  % USER: fig_anchor_bank_selection.pdf
\end{figure}

Canonical half-body selection matters because full-body anchors often include distracting pose or background context that IP-Adapter over-copies into story panels.

\subsection{IP-Adapter on versus off}

\begin{figure}[!h]
  \centering
  \begin{subfigure}{0.48\columnwidth}
    \fbox{\parbox[c][0.75in][c]{0.95\linewidth}{\centering\scriptsize without IP-Adapter\\PLACEHOLDER}}
    \caption{Off}
  \end{subfigure}
  \hfill
  \begin{subfigure}{0.48\columnwidth}
    \fbox{\parbox[c][0.75in][c]{0.95\linewidth}{\centering\scriptsize with IP-Adapter\\PLACEHOLDER}}
    \caption{On}
  \end{subfigure}
  \caption{Single-character story: identity drift without visual conditioning (left) vs.\ anchor-conditioned generation (right).}
  % USER: two-panel comparison from ablation run
\end{figure}

\subsection{Native identity reference failure}

\begin{figure}[!h]
  \centering
  \fbox{\parbox[c][0.85in][c]{0.95\columnwidth}{\centering\scriptsize Identity ref frame\\(character sheet / multi-view)\\PLACEHOLDER}}
  \caption{StoryDiffusion identity reference degenerating into a character sheet, weakening the identity bank.}
  \label{fig:identity_sheet}
  % USER: fig_native_identity_ref_sheet.pdf
\end{figure}

\subsection{DiT supplementary cases}

\begin{figure}[!h]
  \centering
  \begin{subfigure}{0.48\columnwidth}
    \fbox{\parbox[c][0.75in][c]{0.95\linewidth}{\centering\scriptsize DiT best case\\PLACEHOLDER}}
    \caption{Best}
  \end{subfigure}
  \hfill
  \begin{subfigure}{0.48\columnwidth}
    \fbox{\parbox[c][0.75in][c]{0.95\linewidth}{\centering\scriptsize DiT failure\\PLACEHOLDER}}
    \caption{Failure}
  \end{subfigure}
  \caption{Exploratory PixArt-$\alpha$ runs: occasional coherent panel (left) vs.\ identity or artifact failure (right).}
  % USER: optional extra DiT figures
\end{figure}
